\chapter{Conclusion and Future Directions}\label{chapter:conclusion-future-directions}

In this work we have developed an implementation of the \ac{FMM} for the \ac{BEM} formulation of the Helmholtz equation for sound hard exterior scattering problems. We began with a classical collocation point based implementation of the \ac{BEM} made with the \texttt{deal.II} library, and proceed to implement the \ac{BM} formulation of the problem, which mitigates error spikes at certain fictitious frequencies. We built an Octree data structure for the \ac{FMM} that adpatively subdivides the domain into boxes. We added Z-order space filling curves to make tree traversal more efficient. We implemented the translation matrices for the \ac{FMM} and the tree traversal with the upward, downward, and horizontal passes. We implemented the matrix-vector multiplication for the \ac{FMM} which is the crux of the algorithm that lets us solve the sytem without ever having to explicitly assemble the full matrix. We implemented a conversion layer that converts our complex system into the real-equivalent system that allows us to utilize the native \texttt{deal.II} implementation of \ac{GMRES} to iteratively solve the system. Finally, we leveraged OpenMP strategically to parallelize the computation of the translation matrices and the matrix-vector multiplication for shared memory systems.

We validated our implementation against a known analytical solution for plane-wave scattering off of a single sphere. We also demonstrated the effectiveness of the \ac{BM} formulation in mitigating error spikes at fictitious frequencies. Finally, we demonstrated the measurable improvement in the computational order provided by the \ac{FMM} over a traditional direct \ac{BEM} approach.

While our implementation has demonstrated its performance, it stil has some limitations that we must acknowledge. The current implementation is not well suited for higher frequencies. we have implemented the original classical mutlipole expansions uusing regular spherical harmonics. As the frequency goes up, we need an ever smaller final leaf node size to maintain accuracy, which leads to increasing depth and increasing number of translation matrices that need to be stored. This can be addressed by implementing the plane wave expansion of the multipole expansions, which is more efficient for higher frequencies, but breaks down for lower frequencies. They can be used together in a combined approach as demonstrated by \cite{gumerovbroadbandfastmultipole2009}. 
Our \ac{BM} implementation depends on regularization that arises from the use of a single dof per cell and perfectly flat panels. This is not the case for higher order elements, and so we would need to implement a more general regularization approach to support higher order elements. We have also built the solver around purely the sound-hard scatterer problem, and this can be extended to support arbitrary impedance boundary conditions. It would also be interesting to use an iterative solver that natively supports complex numbers and compare it to the convergence of the real-equivalent system we use with \ac{GMRES}. Also, the use of flexible iterative solvere with adjustable preconditioners, where the preconditioning is based on the solution of a coarser \ac{FMM} itself could be interesting to implement. Finally, we have only implemented a shared memory parallelization strategy with OpenMP, but the \ac{FMM} lends itself to distributed memory parallelization as well, and this can be implemented with MPI.
